\section{第二部分实验：真值、攻击与解释}\label{sec:noderes}

\subsection{没有真值，就造一个}

Cora 上没有“正确的解释”，所以要造一个。合成图用随机块模型：3 个社区各 100 节点，$p_{\text{in}}=0.06$、$p_{\text{out}}=0.003$，节点特征是“社区 one-hot + 高斯噪声”。标签由社区决定，于是“哪些边携带标签信号”是已知的。

\textbf{真值}：对每条候选边 $e$，删掉它重算目标类的概率，下降量
\begin{equation}
\Delta_e = p(c\mid G) - p\bigl(c\mid G-e\bigr)
\end{equation}
就是它的真实重要性。这是一个\emph{直接测出来}的量，不依赖任何假设。

有了 $\Delta$，就可以问三个问题：解释器选的边平均 $\Delta$ 有多高？与随机边比如何？与“按 $\Delta$ 排序取 top-$k$”的 oracle 比如何？

\subsection{攻击：FGA 比随机强 7 倍}

\begin{table}[H]
\centering
\caption{Cora 上的边注入攻击（准确率下降）}
\label{tab:attack}
\small
\begin{tabular}{lcccc}
\toprule
预算 $\varepsilon$ & 1\% & 3\% & 5\% & 10\% \\
\midrule
随机加边 & $-0.002$ & $-0.009$ & $-0.012$ & $-0.017$ \\
FGA & $-0.016$ & $-0.046$ & $-0.075$ & $\mathbf{-0.119}$ \\
\bottomrule
\end{tabular}
\end{table}

干净准确率 0.814。FGA 在 10\% 预算下降 11.9 个点，是随机加边的 7 倍，且两条曲线都单调（表~\ref{tab:attack}）。值得注意的是：即使加 528 条边，模型仍保有 69.5\% 的准确率——这个模型的鲁棒性比预期强。

\subsection{解释质量：合成环境里的真值检验}

\begin{figure}[H]
\centering
\includegraphics[width=0.84\textwidth]{b_explain.png}
\caption{左：解释边、oracle 上界与随机边的平均消融重要性。右：逐节点效率散点（越靠近对角线越好）。}
\label{fig:bexplain}
\end{figure}

在 SBM 上，44 个可判定节点中（图~\ref{fig:bexplain}）：解释边平均 $\Delta=0.0673$，随机边 $0.0044$（Wilcoxon 符号秩 $p=5\times10^{-11}$）；\textbf{oracle 效率中位 0.939}——即逐节点“解释边的重要性 / 最优 top-$k$ 的重要性”的中位数为 0.939，解释器达到了“直接按真值排序”的 94\%。候选池内百分位 0.669（随机为 0.5）。

（这里要说清统计量：均值之比是 0.901，中位数之比是 0.939。我们报中位数，因为它不被个别节点拉偏；两个数都支持同一结论。）

还有 1/45 个节点的掩码并列、排不出序。这个比例不高，但它说明“解释”并非总有定义。

\subsection{Cora 上的解释：只能测两件事}

在没有真值的真实数据上，只能测\textbf{充分性}与\textbf{必要性}。

\textbf{充分性}：只保留解释边，目标类概率还剩多少？解释边为原来的 1.31 倍，随机 incident 边为 1.21 倍——解释边更好。

\textbf{必要性}：删掉解释边，目标类概率降多少？解释边降 0.123，随机 incident 边只降 0.019——差了一个数量级。

\textbf{稳定性}：换一个训练种子重新训模型，同一个节点的解释还一样吗？跨 3 个种子的 Jaccard 中位为 0.714，说明解释基本稳定，但确实会随模型变化。

掩码解释器在 Cora 上的可判定率为 88\%（105/120）——候选边中位 48 条时它仍然能给出排序。

\subsection{一个反直觉的发现：攻击不改变解释}

用 FGA 加 5\% 的边把准确率打掉 7.5 个点，然后重新解释同一批节点，比较攻击前后的解释边集合。按\emph{节点对}比较（注意：按边下标比较是错的——攻击后边表重排，5278 条原始边里只有 24 条保持同一下标），\textbf{Jaccard 中位为 1.000}：解释几乎逐条相同。

如何理解？解释器给出的边反映的是“哪些邻居对目标节点的表示重要”这种局部结构，而 FGA 加的是跨社区边——它改变了模型的输出，却没有改变目标节点最重要的那几个邻居。所以解释稳定，但它也\emph{不会告诉你攻击改动了哪里}。

这对可解释性是双刃：稳定是好事（解释不受噪声影响），但一个在模型明显变差时仍然纹丝不动的解释，未必能作为“审计工具”。

\subsection{负控制}

只测“能跑通”的断言等于没测。我们设计三条负控制，要求判据必须失败：

\begin{table}[H]
\centering
\caption{三条负控制}
\label{tab:neg}
\small
\begin{tabular}{llc}
\toprule
桩实现 & 期望 & 实测 \\
\midrule
用 MLP 顶替 GCN & 准确率掉到 $\le0.65$ & 0.502 \\
打乱标签 & 准确率塌到多数类附近 & 0.286（多数类 0.302） \\
用 bottom-$k$ 顶替 top-$k$ 解释 & 消融重要性倍率 $<2$ & $0.09\times$ \\
\bottomrule
\end{tabular}
\end{table}

第三条尤其重要：它证明主判据有“牙”——如果解释器挑的是最不重要的边，结论会立刻塌掉（表~\ref{tab:neg}）。
